% Part: lambda-calculus
% Chapter: church-rosser
% Section: definitions-and-properties

\documentclass[../../../include/open-logic-section]{subfiles}

\begin{document}

\olfileid[pt]{lam}{cr}{dap}

\olsection{Definição e propriedades}

Neste capítulo, introduzimos o conceito da propriedade de Church--Rosser e
algumas propriedades comuns dessa propriedade.

\begin{defn}[Propriedade de Church--Rosser, CR] Diz-se que uma relação
  $\xredone$ sobre termos satisfaz a \emph{propriedade de Church--Rosser} se e
  somente se, sempre que $M \xredone P$ e $M \xredone Q$, existe algum~$N$
  tal que $P \xredone N$ e $Q \xredone N$.
\end{defn}

Podemos considerar o cálculo lambda como um modelo de computação no qual os
termos em forma normal são «valores» e uma relação de redutibilidade sobre
termos constitui as «regras de cálculo». A propriedade de Church--Rosser
afirma que, quando há mais de uma maneira de prosseguir com um cálculo, ainda
assim existe um único valor da expressão.

Como exemplo da álgebra elementar, há mais de uma maneira de calcular
$4 \times (1+2) + 3$. A expressão pode ser reduzida a $4 \times 3+3$ (se
primeiro reduzirmos $1+2$ a~$3$) ou a $4 \times 1+4 \times 2+3$ (se primeiro
reduzirmos $4 \times (1+2)$ usando a distributividade). Contudo, ambas podem
ser reduzidas posteriormente a $12+3$.

Se tomarmos $\xredone$ como a redução~$\beta$, veremos facilmente que uma
consequência da propriedade de Church--Rosser é que, se um termo tem uma forma
normal, ela é única. De fato, suponha que $M$ possa ser reduzido a $P$ e a
$Q$, ambos formas normais. Pela propriedade de Church--Rosser, existe algum
$N$ ao qual tanto $P$ como $Q$ se reduzem. Como, por hipótese, $P$ e $Q$ são
formas normais, a redução de $P$ e $Q$ a $N$ só pode ser a redução trivial;
O.s.\@, $P$, $Q$ e $N$ são idênticos. Isso justifica falarmos em \emph{a}
forma normal de um termo.

Ao considerar o cálculo lambda como um modelo de computação, portanto, a
forma normal de um termo pode ser vista como o «resultado final» da computação
que começa com esse termo. O corolário acima significa que existe, se houver
algum, apenas um resultado final de uma computação, assim como há apenas um
resultado do cálculo de $4 \times (1+2)+3$, a saber,~$15$.

\begin{thm} \ollabel{thm:str}
  Se uma relação $\xredone$ satisfaz a propriedade de Church--Rosser, e
  $\xred$ é a menor relação transitiva que contém $\xredone$, então $\xred$
  também satisfaz a propriedade de Church--Rosser.
\end{thm}

\begin{proof}
  Suponha
  \begin{align*}
    M & \xredone P_1 \xredone \dots \xredone P_m \text{ e}\\
    M & \xredone Q_1 \xredone \dots \xredone Q_n.
  \end{align*}
  Demonstraremos o teorema construindo uma grade $N$ de termos com altura
  $m + 1$ e largura $n + 1$. Usamos $N_{i,j}$ para denotar o termo na
  $i$-ésima linha e na $j$-ésima coluna.
  
  Construímos $N$ de modo que $N_{i,j} \xredone N_{i+1,j}$ e
  $N_{i,j} \xredone N_{i,j+1}$. Ela é definida da seguinte maneira:
  \begin{align*}
    N_{0,0} &= M \\
    N_{i,0} &= P_i && \text{se } 1 \le i \le m \\
    N_{0,j} &= Q_j && \text{se } 1 \le j \le n \\
  \intertext{e, nos demais casos:}
    N_{i,j} &= R
  \end{align*}
  onde $R$ é um termo tal que $N_{i-1,j} \xredone R$ e $N_{i,j-1}
  \xredone R$. Pela propriedade de Church--Rosser de $\xredone$, tal termo
  sempre existe.

  Agora temos $N_{m,0} \xredone \dots \xredone N_{m,n}$ e $N_{0,n}
  \xredone \dots \xredone N_{m,n}$. Observe que $N_{m,0}$ é $P$ e $N_{0,n}$
  é $Q$. O teorema decorre da definição de~$\xred$.
\end{proof}

\end{document}
